For this, consider the complex:
@@M0@@
where @@M1@@ is the symmetric algebra of @@M2@@. This is a resolution of @@M3@@ over @@M4@@ (VI 1.11). The associated complex on @@M5@@, denoted @@M6@@, has @@M7@@-th term @@M8@@ and is therefore a resolution of @@M9@@ by locally free @@M10@@-Modules of finite type; thus we have @@M11@@. Consequently:
@@M12@@
The proposition follows by additivity.

\sgaStableTarget{sga6.e09.sec.2}\subsection*{2. Principal bundles under split tori}

\sgaStableTarget{sga6.e09.para.2-1}2.1. Let @@M13@@ be a noetherian scheme, and @@M14@@ a torus over @@M15@@. Suppose @@M16@@ split; we thus have an isomorphism @@M17@@, where @@M18@@ is a free abelian group of rank @@M19@@, or equivalently @@M20@@. Let @@M21@@ be a torsor under @@M22@@. Then @@M23@@ is affine over @@M24@@, defined by a graded @@M25@@-Algebra @@M26@@ of type @@M27@@; putting @@M28@@, we have that @@M29@@ is an invertible @@M30@@-Module for every @@M31@@, and, for @@M32@@, the homomorphism @@M33@@ defined by multiplication in @@M34@@ is an isomorphism (SGA 3 VIII 4.1). We thus define a multiplicative homomorphism @@M35@@ by associating to @@M36@@ the class of @@M37@@ in @@M38@@.
